% =============================================================================
% Lecture 18: Multi-Agent Coordination
% Agentic Machine Learning Systems (Volume III), Chapter 18
% =============================================================================
% Thesis: splitting a task adds boundaries one agent did not have, so
% delegation earns its place only when it delivers more accepted tasks per unit
% of cost, or sooner in a way the objective values, than a capable single
% controller with the same budget.
% Hook and synthesis: Anthropic's multi-agent research system (War Story 18.1;
% Hadfield et al. 2025): about 15 times the tokens of chat, 90.2 percent better
% on an internal evaluation. Running case: the repair agent's authentication
% migration, migration_task_01, split between worker_interface_01 and
% worker_caller_01 (MA §18.1 to §18.9).
%
% Build: make deck V=3 D=18_multi_agent   (figures: images/make_figs.py)
% Plan:  lecture_map.yml in this folder.
% Numbers: images/calcs.py prints every derived value; sources named there.
% =============================================================================

\documentclass[aspectratio=169, 12pt]{beamer}
\usepackage{../../assets/beamerthememlsys}
\mlsyssetup{
  volume={Volume III},
  chapter={Chapter 18},
  chaptertitle={Multi-Agent Coordination},
  layer=scale,
  logo={../../assets/img/logo-mlsysbook.png},
  instlogo={../../assets/img/logo-harvard.png},
}
\usepackage[T1]{fontenc}
\usepackage[scaled=0.9]{helvet}
\usepackage{courier}
% Code and real tool output: Inconsolata is narrower than Courier.
\usepackage[varqu,scaled=0.95]{zi4}
\renewcommand{\familydefault}{\sfdefault}
\usepackage{booktabs,amsmath}
\graphicspath{{images/}}
\setcounter{mlsystotalsections}{5}

\title{Multi-Agent Coordination}
\author{Vijay Janapa Reddi}
\institute{Harvard University}
\date{}

% A large answer choice for predict and concept-question slides
\newcommand{\choice}[2]{\textbf{\textcolor{mlsysaccent}{#1}}\enspace #2}
% A choice that is highlighted from overlay #1 on (the answer)
\newcommand{\answerat}[3]{\alt<#1->{\colorbox{mlsysaccent!15}{\choice{#2}{#3}}}{\choice{#2}{#3}}}
% A goal or takeaway line with its thumbnail
\newcommand{\thumbline}[2]{%
  \begin{minipage}[c]{2.6cm}\centering\includegraphics[width=2.5cm,height=0.95cm,keepaspectratio]{#1}\end{minipage}%
  \hspace{0.3cm}\begin{minipage}[c]{0.74\textwidth}\small\raggedright #2\end{minipage}}
% An identifier from the running case, in code type
\newcommand{\op}[1]{\texttt{#1}}
% An exact excerpt of real tool output, a config, code, or a document
\newenvironment{logblock}[1][7.2]{\par\ttfamily\fontsize{#1}{8.6}\selectfont\raggedright\setlength{\parskip}{0pt}}{\par}
% Code indentation of n spaces in the current monospace font
\newcommand{\ind}[1]{\hspace*{#1\fontcharwd\font`0}}
% Golden-thread tag: one line on the goals slide naming the H·S·A term this lecture acts on
\newcommand{\hsathread}[1]{{\scriptsize\textbf{\textcolor{mlsysaccent}{H·S·A thread:}} #1\par}}
% Colors used in the figures, for text beside them
\definecolor{figviolet}{HTML}{7E317B}
\definecolor{figred}{HTML}{CB202D}
\definecolor{figgreen}{HTML}{008F45}
\definecolor{figorange}{HTML}{E67817}

\begin{document}

% =============================================================================
% TITLE AND HOOK
% =============================================================================
\mlsysnote{[0.5 min] [CORE] Read the thesis aloud once. Lecture 17 trained one agent from its own trajectories. Today the question changes from how good one agent is to whether several of them beat one. Notes total 71 minutes; the core lecture is 59 minutes. FLEX: notes tagged [OPTIONAL] can be cut. Cut first the scalability-law plot, the group-chat code, the Erlang supervisor, the shared-write plot, the Shuttle computers, the credit window, and the open questions. Every concept the chapter owns stays on a core slide.}
\mlsystitle{Multi-Agent Coordination}{Splitting a task adds boundaries one agent did not have, so delegation pays only when it beats one equally funded controller on cost per accepted task or on a deadline.}{cover_multi_agent_blueprint.png}

\begin{frame}{Predict: did 15 times the tokens pay?}
  \mlsysactivity{Predict}{3 min}
  \begin{columns}[c]
    \begin{column}{0.56\textwidth}
      \centering\safeimg[width=\textwidth,height=4.0cm,keepaspectratio]{research_arch.png}\\[0.1cm]
      \begin{logblock}[7]
        In our data, agents typically use about 4\texttimes{} more tokens than chat interactions, and multi-agent systems use about 15\texttimes{} more tokens than chats.
      \end{logblock}
    \end{column}
    \begin{column}{0.41\textwidth}
      \small The multi-agent system scored \textbf{90.2\%} higher than one agent on the team's own evaluation.\\[0.2cm]
      Per accepted task, is it cheaper?\\[0.15cm]
      \choice{A}{yes, always}\\[0.05cm]
      \choice{B}{no, never}\\[0.05cm]
      \answerat{2}{C}{only on tasks one agent usually fails}\\[0.15cm]
      \uncover<2->{\footnotesize Break-even needs one agent to succeed at most \textbf{26.7\%} of the time.}
    \end{column}
  \end{columns}
  \mlsyscredit{Hadfield et al., ``How we built our multi-agent research system,'' Anthropic Engineering, 13 June 2025 (figure and excerpt). Multi-Agent Coordination, War Story 18.1 and Napkin Math 18.6.}
  \note{[3.5 min] [CORE] NARRATE: a real system in its builders' own figure. A lead agent plans a query, spawns parallel search subagents with their own contexts, and writes the report from what they return. The excerpt is their measurement: an agent uses about 4 times the tokens of a chat, and the multi-agent system about 15 times. They report that it beat a single agent by 90.2 percent on their internal research evaluation. ENGAGE: one minute; vote. Step 2: C. The delegated design spends 15 divided by 4, or 3.75 times the tokens of one agent per attempt. Cost per accepted task is attempt cost over acceptance, so delegation is cheaper only if it multiplies acceptance by more than 3.75. Acceptance cannot exceed 1, so that is possible only where one agent succeeds at most 26.7 percent of the time (Napkin Math 18.6, derived). Even if the 90.2 percent gain were a 1.902 times acceptance multiple, cost per accepted task would be 1.97 times the single agent's. WARN: the 90.2 percent is a score on one team's evaluation and measures no acceptance rate; the chapter treats it as one reported system. LINK: cost per accepted task is Lecture 1's relation; today asks when splitting improves it. FLEX: [CORE].}
\end{frame}

\begin{frame}{Recall: how many rewarded rollouts were failures?}
  \mlsysactivity{Retrieval}{1.5 min}
  \begin{columns}[c]
    \begin{column}{0.56\textwidth}
      {\Large\bfseries Close your notes.}\\[0.2cm]
      {\normalsize 8,192 sampled attempts (rollouts), 10\% truly pass. The verifier accepts every pass and 2.5\% of failures. What share of rewards go to failures?}\\[0.25cm]
      \uncover<2->{\small \textbf{\textcolor{mlsysaccent}{18.4\%}}: 184.32 false accepts against 819.20 true ones.\\[0.08cm]
      A small error rate becomes a large share when passes are rare.}
    \end{column}
    \begin{column}{0.40\textwidth}
      \centering\safeimg[width=\textwidth,height=1.3cm,keepaspectratio]{thumb_recall_bars.pdf}\\[0.25cm]
      \safeimg[width=0.4\textwidth,height=1.7cm,keepaspectratio]{hsa_locator_hsa.pdf}\\[0.02cm]
      {\footnotesize All three exposures lit: horizon, state, authority}
    \end{column}
  \end{columns}
  \mlsyscredit{Reinforcement Learning from Verifiable Rewards, Napkin Math 17.1 (\S17.3; illustrative). Locator: Multi-Agent Coordination, \S18.1. Derived in images/calcs.py.}
  \note{[1.5 min] [CORE] Retrieval from Lecture 17. Give 45 seconds. Step 2: 8,192 times 0.10 is 819.2 true accepts; 8,192 times 0.90 times 0.025 is 184.32 false accepts; 184.32 over 1,003.52 is 18.4 percent. A checker that is wrong on 2.5 percent of failures still pays 18.4 percent of its rewards to failures. Today that arithmetic returns with model reviewers: what does agreement among several fallible checkers certify? NARRATE the locator: every handoff between agents multiplies all three exposures. Each child runs its own horizon of turns, carries its own state, and holds delegated authority, so the runtime must keep closure across every boundary it adds. FLEX: [CORE].}
\end{frame}

\begin{frame}{When does splitting a task across agents pay?}
  \thumbline{thumb_screen.pdf}{Estimate the useful worker count from the serial fraction and the per-worker overhead.}\\[0.06cm]
  \thumbline{thumb_structures_step3.pdf}{Separate dependencies, messages, and ownership, and bound a delegation tree before it runs.}\\[0.06cm]
  \thumbline{thumb_merge_step3.pdf}{Explain why a clean merge can still be invalid, and publish only a checked candidate.}\\[0.06cm]
  \thumbline{thumb_floor_step2.pdf}{Compute what agreement among correlated reviewers can certify.}\\[0.06cm]
  \thumbline{thumb_cancel_step3.pdf}{Bound in-flight work, authority, and spending after a split.}\\[0.06cm]
  {\footnotesize\textcolor{midgray}{Running case: the repair agent migrates an authentication module, split between two workers.}}\\[0.04cm]
  \hsathread{multiplies all three exposures}
  \note{[1 min] [CORE] The driving question is the title. Five goals, one per segment; each thumbnail is the shape of the figure that teaches it. The running case is the repair agent from Lecture 1 on a wider job: migrate an authentication module, which changes an interface and every caller of it. One controller could keep the whole task. The split tempts because the work is wide: an interface worker and a caller worker, each with a smaller context. Today we price that split.}
\end{frame}

% =============================================================================
% SECTION: SPLITTING
% =============================================================================
\mlsysnote{[0.5 min] [CORE] First: what does a split add that one controller lacked, and when can it pay at all?}
\mlsysdivider{What does a split add, and when can it pay?}{Multi-Agent Coordination, \S18.2 and \S18.9}
\section{Splitting}

\begin{frame}{A delegated worker owns a whole child trajectory}
  \centering
  \mlsyssteps[width=0.94\textwidth,height=5.4cm,keepaspectratio]{split_step}{2}
  \mlsyscredit{Multi-Agent Coordination, \S18.2 and Table 1.}
  \note{[2 min] [CORE] Step 1: the capable baseline. One controller already runs searches, tests, and candidate samples in parallel, keeps its context lean with retrieval and compaction, and holds a scoped grant. Parallel calls inside one trajectory are not delegation; Lecture 3's self-consistency votes are the cheapest ensemble and also the most correlated, because every sample shares the model, prompt, and input. Step 2: delegation creates child trajectories, each with its own context, model-and-tool loop, task contract, and lifecycle. What it adds is the red boundaries: each must carry explicitly what one trajectory carried implicitly, the inputs a child may rely on, what it may do and spend, and what it must return. ENGAGE: which of these does delegation buy that the baseline lacks? (Expected: separate working contexts, task-scoped authority, independent multi-turn ownership; the baseline already has parallelism.) WARN: if both workers need the full repository history, separate contexts duplicate prompt state instead of isolating it. FLEX: [CORE].}
\end{frame}

\begin{frame}{The serial part caps speedup; overhead grows per worker}
  \begin{columns}[c]
    \begin{column}{0.26\textwidth}
      \centering\safeimg[width=\textwidth,height=4.0cm,keepaspectratio]{amdahl.jpg}\\[0.04cm]
      {\footnotesize Gene Amdahl, 2008}
    \end{column}
    \begin{column}{0.71\textwidth}
      \centering
      {\large $\text{Speedup}(N)=\dfrac{1}{(1-p)+\dfrac{p}{N}+c(N)}$}\\[0.35cm]
      \small\raggedright
      $p$: the fraction of the work that divides among $N$ workers\\[0.08cm]
      $c(N)$: \textbf{coordination tax}, the time spent on handoffs, as a fraction of the serial time\\[0.08cm]
      With $c=0$: \textbf{Amdahl's law}, capped at $1/(1-p)$\\[0.12cm]
      \uncover<2->{\textcolor{mlsysaccent}{\textbf{Any gain needs $c(N) < p\,(1-1/N)$.}}}
    \end{column}
  \end{columns}
  \mlsyscredit{Photo: Perry Kivolowitz, CC BY 3.0, Wikimedia Commons. Multi-Agent Coordination, equation 1 and Definition 18.1; Amdahl 1967.}
  \note{[2 min] [CORE] Gene Amdahl argued in 1967 that the serial housekeeping of a program limits what any number of processors can deliver; his paper has no equation, and later authors wrote the formula. For agents the housekeeping is planning, integration, and final assessment, which no added worker absorbs. The chapter adds the coordination tax $c(N)$: preparing and validating envelopes, delivering context, scheduling children, collecting receipts, and reconciling results, normalized by the serial time. Step 2: a split beats the serial reference only if the tax is smaller than the time the parallel part saves. ENGAGE: with $p = 0.8$, what is the largest possible speedup with free coordination? (Expected: 5 times, at infinitely many workers.) WARN: the serial reference here is the partition run in sequence; the real baseline already overlaps tools, so a production decision benchmarks the capable controller instead. FLEX: [CORE].}
\end{frame}

\begin{frame}{Six workers minimize elapsed time here, at 2.3 times}
  \centering
  \safeimg[width=0.92\textwidth,height=5.1cm,keepaspectratio]{screen.pdf}\\[0.05cm]
  \small Overhead $c(N) = 0.02\,(N-1)$: eight workers take 0.440 of the serial time, six take 0.433.
  \mlsyscredit{Multi-Agent Coordination, Napkin Math 18.1 (\S18.2.2; hypothetical workload). Derived in images/calcs.py.}
  \note{[1.5 min] [CORE] Napkin Math 18.1. The work divides with $p = 0.80$ and each added worker costs 0.02 of the serial time. Four workers need 0.460 of the serial time; the minimum is at 6 workers, 0.433, a 2.3 times speedup; eight workers rise to 0.440, because the shrinking parallel term no longer pays for the added overhead. Halving the overhead to 0.01 moves the optimum to 9 and the speedup only to 2.7 times. The optimum sits near $\sqrt{p/c}$, 6.3 here. ENGAGE: which lever moves the ceiling? (Expected: only the serial fraction; no worker count beats 5 times.) WARN: students read the gentle curve as harmless; past the peak, every worker costs tokens and buys nothing. FLEX: [CORE].}
\end{frame}

\begin{frame}{Estimate: halve the overhead or the serial work?}
  \mlsysactivity{Estimate}{2.5 min}
  \begin{columns}[c]
    \begin{column}{0.44\textwidth}
      \begin{mlsyscard}{mlsysaccent}
        \small\textbf{Given (Checkpoint 18.1)}\\[0.05cm]
        $p = 0.9$, $c(N) = 0.02\,(N-1)$\\[0.03cm]
        Optimum near $\sqrt{p/c}$; ceiling $1/(1-p)$\\[0.03cm]
        Fund one: halve $c$, or halve $1-p$.
      \end{mlsyscard}
    \end{column}
    \begin{column}{0.53\textwidth}
      \small
      \uncover<2->{Now: $\sqrt{45} = 6.71$; optimum \textbf{7}, \textbf{2.87}$\times$; ceiling 10$\times$\\[0.15cm]}
      \uncover<3->{Halve $c$: 9 workers, \textbf{3.57}$\times$, ceiling still 10$\times$\\[0.1cm]
      Halve serial: 7 workers, \textbf{3.27}$\times$, ceiling \textbf{20}$\times$\\[0.2cm]
      \textcolor{mlsysaccent}{\textbf{Halving $c$ wins now with 2 more workers; only the serial cut raises the ceiling.}}}
    \end{column}
  \end{columns}
  \mlsyscredit{Multi-Agent Coordination, Checkpoint 18.1 and Exercise 1 (hypothetical inputs). Derived in images/calcs.py.}
  \note{[2.5 min] [CORE] ENGAGE: two minutes in pairs with the given card. Step 2: $\sqrt{0.9/0.02} = 6.71$, and the optimal integer count is 7, at 2.87 times; the ceiling is 1/(1-0.9), 10 times. Step 3: halving the overhead to 0.01 gives 9 workers at 3.57 times; halving the serial work to $p = 0.95$ gives 7 workers at 3.27 times and doubles the ceiling to 20 times. Both numbers come from images/calcs.py. Read the trade with token spend in mind: the overhead cut needs two more workers, each with its own tokens, to reach its speedup, while the serial cut reaches 3.27 times with the same seven. WARN: the chapter's general lesson is that the serial fraction sets the ceiling and the overhead sets the useful fleet size; which proposal to fund also depends on the deadline and the token price. FLEX: [CORE].}
\end{frame}

\begin{frame}{Past its peak, each added worker lowers capacity}
  \begin{columns}[c]
    \begin{column}{0.56\textwidth}
      \centering\safeimg[width=\textwidth,height=5.3cm,keepaspectratio]{usl.pdf}
    \end{column}
    \begin{column}{0.41\textwidth}
      \small\textbf{Universal scalability law:}\\[0.05cm]
      $\dfrac{N}{1+\sigma(N-1)+\kappa N(N-1)}$\\[0.15cm]
      $\sigma$: work that serializes\\
      $\kappa$: pairwise reconciliation\\[0.15cm]
      Repair: $\sigma = 0.20$, $\kappa = 0.02$\\
      Research: $\sigma = 0.10$, $\kappa = 0.002$
    \end{column}
  \end{columns}
  \mlsyscredit{Multi-Agent Coordination, Theorem 18.1 and Lighthouse 18.1 (stylized coefficients); Gunther 2008. Derived in images/calcs.py.}
  \note{[2 min] [OPTIONAL] Neil Gunther's universal scalability law separates two costs with different fixes. Contention, $\sigma$, is work that serializes: planning, final assessment, and queueing at the supervisor or the integration queue. Coherency, $\kappa$, is pairwise reconciliation across the $N(N-1)$ ordered worker pairs. With $\kappa = 0$ it is Amdahl's law; with $\kappa > 0$ capacity peaks at $\sqrt{(1-\sigma)/\kappa}$ and then falls. Violet: the repair partition peaks at 6 workers, 2.3 times, and from 41 workers on finishes later than one worker (derived: 0.981 at 41). Green: the research agent's search paths barely talk to each other, so it peaks at 21 subagents and 5.5 times, a 3.5 times wider fan-out. ENGAGE: what would you change to move the violet peak right? (Expected: cut coherency, for example by landing the interface before callers start.) WARN: the coefficients are stylized; in practice teams fit them from trials at a few counts such as 1, 2, and 4, and treat the fit as a hypothesis outside that range. FLEX: [OPTIONAL].}
\end{frame}

\begin{frame}{A 3.75 times token multiple pays only on hard tasks}
  \begin{columns}[c]
    \begin{column}{0.53\textwidth}
      \centering\safeimg[width=\textwidth,height=5.2cm,keepaspectratio]{breakeven.pdf}
    \end{column}
    \begin{column}{0.44\textwidth}
      \small Cost per accepted task, delegated over single:\\[0.05cm]
      {\normalsize $m \cdot \dfrac{p_{\text{single}}}{p_{\text{multi}}}$}, \quad $m = 15/4 = 3.75$\\[0.2cm]
      \begin{logblock}[7]
        For economic viability, multi-agent systems require tasks where the value of the task is high enough to pay for the increased performance.
      \end{logblock}
    \end{column}
  \end{columns}
  \mlsyscredit{Multi-Agent Coordination, Napkin Math 18.6 (\S18.9.2). Excerpt: Hadfield et al., Anthropic Engineering, 2025. Derived in images/calcs.py.}
  \note{[2.5 min] [CORE] Napkin Math 18.6 answers the hook. With cost proportional to tokens, Lecture 1's cost per accepted task gives the delegated-to-single ratio $m\,p_{\text{single}}/p_{\text{multi}}$, where $m$ is the token multiple per attempt, 15 over 4, 3.75. Break-even needs delegated acceptance 3.75 times the single agent's. The line is that requirement; it hits acceptance 1 at a single-agent rate of 26.7 percent, and to the right no acceptance rate can pay back the tokens. The builders say the same thing in their own words in the excerpt. ENGAGE: a single agent already accepts 60 percent of migrations. Can delegation win on cost? (Expected: no; it needs 225 percent acceptance. It can only be justified as buying time or coverage that the objective values.) WARN: the 4 and 15 times multiples come from one reported system; measure your own. LINK: Lecture 19 develops the cohort form of this account. FLEX: [CORE].}
\end{frame}

\begin{frame}{A matched baseline gives one controller the same budget}
  \centering
  \safeimg[width=0.94\textwidth,height=4.6cm,keepaspectratio]{matched.pdf}\\[0.1cm]
  \small Against a matched single-controller baseline, read accepted rate, elapsed time, and cost per accepted task together.
  \mlsyscredit{Multi-Agent Coordination, \S18.9.1 and \S18.9.2.}
  \note{[2 min] [CORE] The matched single-controller baseline. A trial that lets the supervisor run tools in parallel while restricting the single controller to serial calls measures that handicap and says nothing about delegation. Lecture 13's matched-budget evaluation fixes the task sample, criteria, and resource envelope; for delegation it must also fix starting fixtures, evidence access, and identical concurrency limits on model calls and tools. Score the integrated migration end to end: an interface patch that passes its own tests still fails the task if the caller patch breaks. Keep every initiated task in the denominator, including timeouts and unknown verdicts; a task still running at the cutoff enters latency as a lower bound. ENGAGE: a team reports median latency over accepted tasks only. What does that hide? (Expected: every slow failure; it flatters both designs.) WARN: if the delegated design also changes the model or prompts, the trial tests a bundle; isolate delegation with an ablation. FLEX: [CORE].}
\end{frame}

% =============================================================================
% SECTION: STRUCTURE
% =============================================================================
\mlsysnote{[0.5 min] [CORE] Land the segment aloud: the serial fraction sets the ceiling, overhead sets the useful fleet, and only a matched trial decides. Now: what structures does a split create?}
\mlsysdivider{Which structures must the runtime keep separate?}{Multi-Agent Coordination, \S18.3 and \S18.4}
\section{Structure}

\begin{frame}{Dependencies, messages, and ownership are three graphs}
  \centering
  \mlsyssteps[width=0.94\textwidth,height=5.4cm,keepaspectratio]{structures_step}{3}
  \mlsyscredit{Multi-Agent Coordination, \S18.3 (S supervisor, I interface worker, C caller worker, T test container).}
  \note{[2.5 min] [CORE] Three structures, one migration. Step 1: the dependency graph records causal precedence. The caller's patch depends on the interface the other worker is changing, and the test run depends on the caller patch. Step 2: the communication topology records who sends messages. Here both workers talk only to the supervisor; they never message each other. Step 3: the supervision tree records lifecycle ownership: the supervisor owns both workers, and the caller owns its test container. ENGAGE: I and C never talk. Can they still produce patches that break each other? (Expected: yes; the dependency edge exists whatever the message routes are.) WARN: a runtime that treats the message graph as the dependency graph will publish a caller patch built on a stale interface; one that treats it as the ownership tree leaves orphaned containers after a cancel. FLEX: [CORE].}
\end{frame}

\begin{frame}{A group chat sends every message to every other agent}
  \begin{columns}[c]
    \begin{column}{0.52\textwidth}
      {\footnotesize\textbf{AutoGen 0.2, GroupChatManager.run\_chat (excerpt)}}\par\vspace{0.05cm}
      \begin{logblock}[7]
        for i in range(groupchat.max\_round):\\
        \ind{4}self.\_last\_speaker = speaker\\
        \ind{4}groupchat.append(message, speaker)\\
        \ind{4}\# broadcast the message to all agents\\
        \ind{4}\# except the speaker\\
        \ind{4}for agent in groupchat.agents:\\
        \ind{8}if agent != speaker:\\
        \ind{12}self.send(message, agent,\\
        \ind{14}request\_reply=False, silent=True)
      \end{logblock}
    \end{column}
    \begin{column}{0.45\textwidth}
      \centering\safeimg[width=\textwidth,height=4.8cm,keepaspectratio]{mesh.pdf}
    \end{column}
  \end{columns}
  \mlsyscredit{microsoft/autogen, branch 0.2 (commit c631b34), autogen/agentchat/groupchat.py, lines 1158 to 1165. Delivery counts: Multi-Agent Coordination, \S18.3.1.}
  \note{[1.5 min] [OPTIONAL] Real framework code: the loop at the heart of AutoGen's group chat, a widely used multi-agent framework. Every time one agent speaks, the manager appends the message to the shared history and sends it to every other agent. When each of $N$ agents speaks once, that is $N(N-1)$ deliveries, against $2N$ for a supervisor that sends one envelope and receives one receipt per worker: 56 against 16 at eight agents, $(N-1)/2$ times more. Each recipient appends the message to its prompt history, so the tokens each model ingests grow with $N^2$ per round, and Lecture 2's append-only resend compounds that across rounds. ENGAGE: why might a team still choose a group chat? (Expected: open-ended discussion where every agent needs every message; the cost is explicit.) WARN: transmission count is not ingested tokens; verbose receipts or frequent polling can exhaust a supervisor's context too. FLEX: [OPTIONAL].}
\end{frame}

\begin{frame}{Diagnose: which node should the team speed up?}
  \mlsysactivity{Diagnose}{3 min}
  \begin{columns}[c]
    \begin{column}{0.62\textwidth}
      \centering\mlsyssteps[width=\textwidth,height=4.6cm,keepaspectratio]{gantt_step}{2}
    \end{column}
    \begin{column}{0.35\textwidth}
      \small Completion time $\ge \max\!\left(\dfrac{W}{N},\ L\right)$\\[0.05cm]
      $W$: total slot work, 275 s\\
      $L$: longest dependency chain\\[0.15cm]
      Cut v6 by 20 s, or v5 by 20 s?\\[0.1cm]
      \uncover<2->{\textbf{v5}: 174.5 to \textbf{154.5} s.\\ v6: no change; 26.5 s of slack.}
    \end{column}
  \end{columns}
  \mlsyscredit{Multi-Agent Coordination, Napkin Math 18.2 and Checkpoint 18.1 (illustrative durations). Derived in images/calcs.py.}
  \note{[2.5 min] [CORE] The slot-pool bound. With $N$ interchangeable slots, no schedule finishes before total slot work divided by slots, nor before the longest dependency chain. Napkin Math 18.2: seven nodes, 275 s of work, 4 slots, so the work bound is 68.75 s; the critical path v1, v2 or v3, v5, v7 is 174.5 s, which the schedule attains with at most 3 slots busy. That is a 1.58 times speedup over running the nodes in sequence, and a fifth slot buys nothing. ENGAGE: three minutes in pairs: cut v6 by 20 s, or v5 by 20 s? Step 2: v5. The v4 to v6 branch finishes 26.5 s before v7 can start, so cutting v6 changes nothing; v5 sits on the path, so completion drops to 154.5 s (derived). Cutting v2 alone also buys nothing, because v3 ties it. WARN: this bound uses assumed durations; Lecture 14 measures the critical path from recorded spans, and retries and queueing stretch real runs beyond either bound. FLEX: [CORE].}
\end{frame}

\begin{frame}{The runtime decides how many children may run}
  \begin{columns}[c]
    \begin{column}{0.44\textwidth}
      \begin{logblock}[7.2]
        Early agents made errors like spawning 50 subagents for simple queries, scouring the web endlessly for nonexistent sources, and distracting each other with excessive updates.
      \end{logblock}
      \vspace{0.2cm}
      \small Their fix was a prompt rule: one agent for a simple fact, more for open-ended work.\\[0.15cm]
      \textcolor{mlsysaccent}{\textbf{Only an admission gate enforces the bound.}}
    \end{column}
    \begin{column}{0.53\textwidth}
      \centering\safeimg[width=\textwidth,height=5.0cm,keepaspectratio]{bounds.pdf}
    \end{column}
  \end{columns}
  \mlsyscredit{Excerpt: Hadfield et al., Anthropic Engineering, 2025. Multi-Agent Coordination, War Story 18.1 and \S18.3.3.}
  \note{[2.5 min] [CORE] Spawning as admission. In early versions of the research system the lead agent decided its own fan-out, and it spawned 50 subagents for simple queries. The team's fix was prompt guidance scaled to query complexity, and a lead that waits for each batch of subagents before the next step. The chapter's lesson: a model may propose a child, but the runtime admits it under four separate bounds. Breadth caps children per parent; concurrent slots cap how many run at once; depth caps delegation layers; a deadline caps time. A shallow tree can still exhaust resources through breadth, and a concurrency cap still allows an unbounded sequence. Last, Lecture 10's bounded-liability admission: reserve each child's maximum charge against the shared budget, atomically, before it starts. ENGAGE: each child gets a budget smaller than the parent's remaining funds. Why is that not enough? (Expected: concurrent siblings can each claim the same remaining balance; the last segment computes it.) FLEX: [CORE].}
\end{frame}

\begin{frame}{A supervisor restarts children under a declared policy}
  \begin{columns}[c]
    \begin{column}{0.50\textwidth}
      {\footnotesize\textbf{Erlang/OTP supervisor (documentation example)}}\par\vspace{0.05cm}
      \begin{logblock}[7]
        init(\_Args) ->\\
        \ind{4}SupFlags = \#\{strategy => one\_for\_one,\\
        \ind{17}intensity => 1, period => 5\},\\
        \ind{4}ChildSpecs = [\#\{id => ch3,\\
        \ind{18}start => \{ch3, start\_link, []\},\\
        \ind{18}restart => permanent, [...]\}],\\
        \ind{4}\{ok, \{SupFlags, ChildSpecs\}\}.
      \end{logblock}
      \vspace{0.15cm}
      \begin{minipage}[c]{0.32\linewidth}\centering\safeimg[width=\linewidth,height=2.0cm,keepaspectratio]{armstrong.jpg}\end{minipage}\hfill
      \begin{minipage}[c]{0.64\linewidth}\footnotesize\raggedright Joe Armstrong, a designer of Erlang, 2009. Too many restarts in a period: the supervisor stops and escalates.\end{minipage}
    \end{column}
    \begin{column}{0.47\textwidth}
      \centering\safeimg[width=\textwidth,height=5.2cm,keepaspectratio]{blast.pdf}
    \end{column}
  \end{columns}
  \mlsyscredit{Erlang/OTP 29 documentation, supervisor principles page (erlang.org/doc/system/sup\_princ.html), excerpt. Photo: Yareeh, CC BY 2.0, Wikimedia Commons. Multi-Agent Coordination, \S18.3.4.}
  \note{[1.5 min] [OPTIONAL] Real code from the Erlang/OTP documentation. Erlang came out of telephone-switch research at Ericsson; Joe Armstrong, pictured, was one of its designers, and his thesis argued that a failing process should crash cleanly while a supervisor restores known-good state. The supervisor declares a strategy and a restart intensity. Under one\_for\_one, only the failed child restarts; fit for disjoint searches. Under one\_for\_all, every sibling in the group restarts; fit when siblings are coupled, as when the interface worker restarts with an altered contract and the caller's partial work becomes invalid. Intensity 1 in period 5 means a second restart within five seconds stops the group and escalates to the parent. ENGAGE: why should an agent restart begin from recorded inputs instead of the failed worker's context? (Expected: a confused context compounds its errors; a restart is a freshly admitted attempt with a fresh reservation.) FLEX: [OPTIONAL].}
\end{frame}

\begin{frame}{An envelope binds the child; a receipt separates claims}
  \begin{columns}[T]
    \begin{column}{0.48\textwidth}
      {\footnotesize\textbf{Task envelope (Listing 1, condensed)}}\par\vspace{0.05cm}
      \begin{logblock}[7]
        "task\_id": "worker\_caller\_01",\\
        "parent\_task\_id": "migration\_task\_01",\\
        "inputs": \{"base\_ref": "repo\_base\_01"\},\\
        "authority": \{"grant\_ref":\\
        \ind{2}"caller\_private\_patch\_grant\_01"\},\\
        "budget": \{"reservation\_ref": "caller\_reservation\_01"\},\\
        "completion": \{"required\_evidence": [...]\},\\
        "return\_limits": \{"receipt\_bytes": 8192,\\
        \ind{2}"optional\_summary\_tokens": 256\}
      \end{logblock}
    \end{column}
    \begin{column}{0.49\textwidth}
      {\footnotesize\textbf{Receipt (Listing 2, condensed)}}\par\vspace{0.05cm}
      \begin{logblock}[7]
        "return\_status": "complete",\\
        \textcolor{figgreen}{"runtime": \{}\\
        \ind{2}\textcolor{figgreen}{"candidate\_ref": "caller\_candidate\_01",}\\
        \ind{2}\textcolor{figgreen}{"write\_set\_ref": "caller\_writes\_01",}\\
        \ind{2}\textcolor{figgreen}{"check\_report\_ref": "caller\_checks\_01",}\\
        \ind{2}\textcolor{figgreen}{"settlement\_ref": "caller\_settlement\_01"\},}\\
        \textcolor{figred}{"worker\_claims": \{"summary":}\\
        \ind{2}\textcolor{figred}{"Updated the designated callers."\}}
      \end{logblock}
    \end{column}
  \end{columns}
  \vspace{0.15cm}
  \centering\small References resolve in runtime records, so a worker cannot widen its grant or budget. \textcolor{figgreen}{Green}: runtime evidence. \textcolor{figred}{Red}: the worker's own words.
  \mlsyscredit{Multi-Agent Coordination, Listings 1 and 2 (\S18.4); fields omitted for space.}
  \note{[2.5 min] [CORE] Typed envelopes and receipts. A prose message like ``update the callers'' leaves open which interface version, which files, what budget, and what counts as done; whatever the worker guesses becomes the contract. Left: the envelope points to records the runtime resolves before dispatch: a pinned repository snapshot, a grant that lets the worker edit caller files in a scratch workspace and withholds publication, and a reservation. The task and dispatch ids tie it to Lecture 11's durable intent. Return limits cap the receipt at 8,192 bytes so bulky output stays in storage. Right: the receipt. Green fields are written by runtime services: the actual diff, the observed dependencies, the exact check run, the usage and settlement records. Red is the worker's summary, kept but unverified. If required evidence is missing, the status becomes evidence\_unavailable, never complete. ENGAGE: Napkin Math 18.3: six children get either a 48,000-token repository slice or a 350-token reference plus 3,600 tokens of reads. Which is cheaper? (Expected: the reference, 23,700 against 288,000 tokens, 12.2 times less, while reads stay under 47,650 tokens each.) FLEX: [CORE].}
\end{frame}

% =============================================================================
% SECTION: MERGING
% =============================================================================
\mlsysnote{[0.5 min] [CORE] Land the segment aloud: three graphs kept separate, admission enforced by the runtime, and a typed contract on every handoff. Now the two workers finish. Can their patches merge?}
\mlsysdivider{Why can two clean patches still break the build?}{Multi-Agent Coordination, \S18.5}
\section{Merging}

\begin{frame}{Predict: different files, no merge conflict. Safe?}
  \mlsysactivity{Predict}{2.5 min}
  \centering
  \mlsyssteps[width=0.94\textwidth,height=5.0cm,keepaspectratio]{merge_step}{3}\\[0.1cm]
  \small Both workers start at \op{repo\_base\_01}. The interface lands first. The caller edits other files.
  \mlsyscredit{Multi-Agent Coordination, \S18.5, figure 3, and Table 4 (simplified: foo stands for the interface).}
  \note{[2.5 min] [CORE] ENGAGE: one minute; vote yes or no before step 3. Step 1: both workers branch from repo\_base\_01, where the interface is foo(a). The caller writes calls to foo(a). Step 2: the interface worker finishes first and its change lands as repo\_head\_02: foo now takes two arguments. Step 3: the caller's patch touches different files, so a three-way textual merge succeeds with no conflict. The merged code calls a function that no longer exists. A clean merge can still be invalid. Disjoint files show only that two edits did not collide on lines, never that they compose. The fallacy as true or false: a clean merge means compatible changes. False. WARN: version control has nothing to warn about here; only a dependency record and a check of the merged candidate catch it. FLEX: [CORE].}
\end{frame}

\begin{frame}{Shared writes pass even odds at five workers}
  \begin{columns}[c]
    \begin{column}{0.56\textwidth}
      \centering\safeimg[width=\textwidth,height=5.3cm,keepaspectratio]{birthday.pdf}
    \end{column}
    \begin{column}{0.41\textwidth}
      \small $\Pr(\text{no shared write}) \approx \exp\!\left(-\dfrac{N(N-1)w^2}{2F}\right)$\\[0.1cm]
      $F$ files, $w$ written by each of $N$ workers\\[0.2cm]
      As in the birthday problem, pairs grow as $N^2$.\\[0.15cm]
      \textcolor{mlsysaccent}{\textbf{Give every worker its own workspace.}}
    \end{column}
  \end{columns}
  \mlsyscredit{Multi-Agent Coordination, equation 4 (\S18.5.1; hypothetical repository, uniform writes). Exact products derived in images/calcs.py.}
  \note{[2 min] [OPTIONAL] Why isolate at all? Equation 4 is the birthday bound generalized to write sets. In a repository of 1,000 files where each worker writes 10 at random, two workers collide 9.6 percent of the time; five workers, 63.9 percent; eight, 94.4 percent. Twenty-three people give even odds of a shared birthday because pairs grow with the square of the group, and so do worker pairs. Real writes cluster on manifests, routing tables, and configuration, so uniform writes understate the risk. In one shared directory three hazards follow: a lost update, where the later write silently replaces the earlier; a torn read, where tests run against a half-changed tree; and lock contention on the version-control index. ENGAGE: Checkpoint 18.2: with 2,000 files, how many workers make a shared write more likely than not? (Expected: 6, derived.) WARN: the other face of the same number: two workers touch disjoint files 90.4 percent of the time, so nearly every pair merges cleanly, which is exactly when the last slide's failure hides. FLEX: [OPTIONAL]; the predict slide already makes the isolation argument.}
\end{frame}

\begin{frame}{Publication checks the exact candidate, then compares and swaps}
  \begin{columns}[c]
    \begin{column}{0.50\textwidth}
      {\footnotesize\textbf{git-worktree(1)}}\par\vspace{0.03cm}
      \begin{logblock}[7]
        A git repository can support multiple working trees, allowing you to check out more than one branch at a time.
      \end{logblock}
      \vspace{0.15cm}
      {\footnotesize\textbf{git-update-ref(1)}}\par\vspace{0.03cm}
      \begin{logblock}[7]
        git update-ref refs/heads/master <new-oid> <old-oid> updates the master branch head to <new-oid> only if its current value is <old-oid>.
      \end{logblock}
      \vspace{0.15cm}
      {\footnotesize Private worktrees share one object store; the publish step is a compare-and-swap on the branch.}
    \end{column}
    \begin{column}{0.47\textwidth}
      \centering\mlsyssteps[width=\textwidth,height=5.4cm,keepaspectratio]{gate_step}{4}
    \end{column}
  \end{columns}
  \mlsyscredit{Git documentation, git-scm.com/docs/git-worktree and git-update-ref (retrieved 2026-10-19), excerpts. Multi-Agent Coordination, \S18.5.2 and Listing 3.}
  \note{[3 min] [CORE] Optimistic concurrency control, from Kung and Robinson's 1981 paper: work privately, validate before anything shared changes, then write. Left: the real mechanisms. A linked worktree gives each agent its own working tree, index, and HEAD over one shared object store; Lecture 8's copy-on-write workspace is the other option. update-ref with an old value is an atomic compare-and-swap on the branch. Right, step by step. 1: compare the receipt's dependency record between the worker's base and the current head; an empty record means unknown, so revalidate. 2: build the exact merge candidate. 3: run the sealed checks on that exact candidate hash, behind Lecture 13's grader isolation; tests that passed on the worker's branch say nothing about the merge. 4: publish only if the head still equals the captured value and the grant is still live. If the head moved, start over at step 1, never blind-retry. A timeout on publish means query the log first. ENGAGE: why does a lease fit a live database migration better than this protocol? (Expected: the database effect cannot be discarded on a failed validation; non-mergeable resources need exclusion at the resource, with Lecture 11's fencing tokens.) FLEX: [CORE].}
\end{frame}

% =============================================================================
% SECTION: REVIEWERS
% =============================================================================
\mlsysnote{[0.5 min] [CORE] Land the segment aloud: isolate each worker, revalidate dependencies, check the exact candidate, publish by compare-and-swap. The merge is guarded; is the code correct? Five reviewer agents vote.}
\mlsysdivider{What does agreement among reviewers prove?}{Multi-Agent Coordination, \S18.6}
\section{Reviewers}

\begin{frame}{The Shuttle's four computers voted; a fifth ran other code}
  \begin{columns}[c]
    \begin{column}{0.45\textwidth}
      \centering\safeimg[width=\textwidth,height=3.6cm,keepaspectratio]{shuttle_gpc.jpg}\\[0.04cm]
      {\footnotesize Space Shuttle general purpose computer units}
    \end{column}
    \begin{column}{0.52\textwidth}
      \begin{logblock}[7]
        Four computers simultaneously execute identical code and synchronize their I/O activities.
      \end{logblock}
      \vspace{0.2cm}
      \begin{logblock}[7]
        In the interest of avoiding a generic software failure, NASA kept its development separate from PASS.
      \end{logblock}
      \vspace{0.2cm}
      {\small A fifth computer ran backup flight software written by another contractor. Voting masks independent faults; identical code fails together.}
    \end{column}
  \end{columns}
  \mlsyscredit{Photo: NASA, public domain, Wikimedia Commons. Spector and Gifford, ``The Space Shuttle Primary Computer System,'' CACM 27(9), 1984; Tomayko, \emph{Computers in Spaceflight}, NASA CR-182505, 1988, Box 4-5.}
  \note{[2 min] [OPTIONAL] Majority voting has a long and successful history in hardware. The Space Shuttle flew four identical computers running the same primary software in a redundant set and voting on outputs, which masks a computer that fails from a particle strike or wear, because those faults are physically independent. NASA did not trust voting to catch a software bug: identical code run four times fails four times. So the fifth computer ran a backup flight system written separately by another contractor, Rockwell, to avoid a generic software failure. Condorcet's jury theorem and von Neumann's majority organs make the same promise, and both rest on the same premise: voters fail independently. ENGAGE: five reviewer agents built on one model family read the same patch. Which kind of redundancy is that? (Expected: four identical computers; their errors share a cause.) LINK: in 1986 Knight and Leveson had 27 versions of one program written independently and found coincident failures far more often than independence allowed. FLEX: [OPTIONAL].}
\end{frame}

\begin{frame}{Different language models often make the same mistake}
  \begin{columns}[c]
    \begin{column}{0.55\textwidth}
      \centering\safeimg[width=\textwidth,height=3.8cm,keepaspectratio]{kim.pdf}
    \end{column}
    \begin{column}{0.42\textwidth}
      \small 71 models on 14,042 multiple-choice questions.\\[0.15cm]
      When two models both answer wrong, they pick the \textbf{same} wrong answer about 60\% of the time.\\[0.15cm]
      Larger, more accurate models agree on errors more.
    \end{column}
  \end{columns}
  \mlsyscredit{Kim, Garg, Peng, and Garg, ``Correlated Errors in Large Language Models,'' ICML 2025, arXiv 2506.07962 (Helm leaderboard; random baseline 1 in 3).}
  \note{[1.5 min] [CORE] Real measurements, cited by the chapter. Kim and coauthors took the responses of 71 language models on the 14,042 multiple-choice questions of the Helm leaderboard. For each pair of models, among questions both got wrong, how often did they pick the same wrong answer? With three wrong options chosen at random, one third of the time. Measured, about 60 percent on average. Models from the same provider or with the same base architecture correlate more, and so do larger, more accurate models even across providers. ENGAGE: what does this predict for a panel of five reviewer agents from different vendors? (Expected: their errors still correlate, so a majority vote buys less than the binomial promises.) WARN: this measures agreement on wrong answers in a benchmark, and false-approval rates on code review are a separate quantity; the chapter's model takes its correlation as an input you must measure on your own known-invalid candidates. FLEX: [CORE].}
\end{frame}

\begin{frame}{Shared difficulty puts a floor under majority error}
  \centering
  \mlsyssteps[width=0.92\textwidth,height=4.5cm,keepaspectratio]{floor_step}{2}\\[0.05cm]
  \small Unanimous false approval of five: 0.032\% if independent; \uncover<2->{\textbf{3.92\%} if correlated.}
  \mlsyscredit{Multi-Agent Coordination, Napkin Math 18.4 and figure 4 (\S18.6.1; beta-binomial model, illustrative). Derived in images/calcs.py.}
  \note{[2.5 min] [CORE] Correlated reviewers. Each reviewer falsely approves 20 percent of known-invalid candidates. Step 1: if errors were independent, a majority of five falsely approves 5.79 percent of the time, and 15 reviewers push it to 0.42 percent. Step 2: real candidates differ in difficulty; a subtle bug fools every reviewer. The chapter models each candidate with a propensity drawn from a beta distribution whose mean is 0.20 and whose pairwise correlation is 0.40. Majority false approval at five is then 16.68 percent, barely better than one reviewer, and it stalls at 15.6 percent at 15 reviewers. As the panel grows, the fraction approving converges to the candidate's propensity, so a large majority approves exactly the candidates with propensity above one half: a floor of 15.0 percent. Tripling the panel triples review spend for one point. ENGAGE: why does the unanimous rate jump 122 times? (Expected: correlation moves probability to the extremes; reviewers fail together or not at all.) WARN: Byzantine agreement tolerates fewer than one third faulty participants; correlated errors push the faulty fraction past that bound however many replicas you add. FLEX: [CORE].}
\end{frame}

\begin{frame}{Vote: ten more reviewers, or one protected test?}
  \mlsysactivity{Vote}{2.5 min}
  \begin{columns}[c]
    \begin{column}{0.55\textwidth}
      \small Five reviewers: 16.68\% majority false approval; the floor is 15.0\%.\\[0.2cm]
      \choice{A}{add ten reviewers of the same kind}\\[0.08cm]
      \answerat{2}{B}{add one protected integration test}\\[0.08cm]
      \choice{C}{require unanimity instead of a majority}\\[0.25cm]
      \uncover<2->{\footnotesize Ten more lower it about one point. Unanimity still fails 3.92\% of the time. A test that runs the merged caller against the new interface checks a fact no reviewer supplies.}
    \end{column}
    \begin{column}{0.42\textwidth}
      \uncover<3->{\footnotesize
      \begin{tabular}{@{}P{2.2cm}P{3.4cm}@{}}\toprule
        \textbf{Evidence} & \textbf{What it grants}\\\midrule
        Reviewer approval & a signal; no authority\\[0.03cm]
        Protected test & one property of one candidate hash\\[0.03cm]
        Publication grant & permission to publish, from the control plane\\\bottomrule
      \end{tabular}}
    \end{column}
  \end{columns}
  \mlsyscredit{Multi-Agent Coordination, Checkpoint 18.2, \S18.6.2, and Fallacy 3 (\S18.10). Derived in images/calcs.py.}
  \note{[2.5 min] [CORE] Peer instruction: vote, one minute with a neighbor, vote again. Step 2: B. More reviewers of the same kind move majority false approval from 16.68 toward the 15.0 percent floor: one point for three times the spend. Unanimity at five still falsely approves 3.92 percent of invalid candidates and also rejects more good ones. The five reviewers in the migration all repeat the same false assumption, that the old interface still exists, so they are one effective observation. A protected test that compiles and runs the merged caller against the new interface supplies a fact they cannot fabricate, behind Lecture 13's grader isolation. The fallacy as true or false: agreement among reviewers proves correctness. False. Step 3: three kinds of evidence, three different grants. Even unanimous approval only advances the candidate to the integration gate. WARN: a test is bounded by its coverage, and a worker that can edit the test can pass it; that is why it stays protected. FLEX: [CORE].}
\end{frame}

% =============================================================================
% SECTION: IN FLIGHT
% =============================================================================
\mlsysnote{[0.5 min] [CORE] Land the segment aloud: correlated reviewers are one signal; buy a different check. Last: what bounds the work, authority, and money still in flight after a split?}
\mlsysdivider{What bounds the work still in flight?}{Multi-Agent Coordination, \S18.7 and \S18.8}
\section{In flight}

\begin{frame}{Vote: each worker is on time 90\%. Is the fan-out?}
  \mlsysactivity{Vote}{2.5 min}
  \begin{columns}[c]
    \begin{column}{0.50\textwidth}
      \centering\uncover<2->{\safeimg[width=\textwidth,height=5.0cm,keepaspectratio]{barrier.pdf}}
    \end{column}
    \begin{column}{0.47\textwidth}
      \small Eight workers must all finish by the deadline.\\[0.1cm]
      \uncover<2->{All eight: $0.9^8$ = \textbf{43.05\%}.\\
      Any 6 of 8 interchangeable: \textbf{96.19\%}.\\[0.15cm]}
      {\footnotesize\textbf{A real fan-out service (latency, ms)}}\\[0.03cm]
      \footnotesize
      \begin{tabular}{@{}lrr@{}}\toprule
        & \textbf{median} & \textbf{99th}\\\midrule
        one random leaf & 1 & 10\\
        all leaves & 40 & 140\\\bottomrule
      \end{tabular}
    \end{column}
  \end{columns}
  \mlsyscredit{Dean and Barroso, ``The Tail at Scale,'' CACM 56(2), 2013, Table 1 (a Google service, measured at the root). Multi-Agent Coordination, Napkin Math 18.5. Derived in images/calcs.py.}
  \note{[2.5 min] [CORE] The fallacy as a vote: if each worker usually finishes on time, so does the fan-out. True or false? Step 2: false. An all-results barrier meets the deadline only if every worker does: $0.9^8$, 43.05 percent, and below even odds already at 7 workers. It is Lecture 1's horizon compounding, applied across workers instead of turns. If the contract accepts any 6 of 8 interchangeable results, as for parallel read-only analyses, the deadline odds recover to 96.19 percent. The real data: Dean and Barroso measured a Google service whose root fans out to many leaves. One leaf answers within 10 ms at the 99th percentile; waiting for all leaves takes 140 ms, 14 times longer. They also report that a hedged second request after a 10 ms delay cut the 99.9th-percentile time to fetch 1,000 values from 1,800 to 74 ms for 2 percent more requests. ENGAGE: can a 6-of-8 rule rescue the migration? (Expected: no; both the interface and caller patches are required. Shrink the fan-out or cut the tail.) WARN: hedging a mutating tool call is a retry before settlement, which Lecture 7 forbids unless both attempts share one operation identity the endpoint deduplicates. FLEX: [CORE].}
\end{frame}

\begin{frame}{A credit window caps the queue and its staleness}
  \begin{columns}[c]
    \begin{column}{0.37\textwidth}
      \centering\safeimg[width=\textwidth,height=3.4cm,keepaspectratio]{kanban2.jpg}\\[0.04cm]
      {\footnotesize A kanban board: each column holds at most its limit}
    \end{column}
    \begin{column}{0.60\textwidth}
      \small Producer: 5 diffs/s. Reviewer: 500 tokens at 40 tokens/s, \textbf{0.08}/s.\\[0.1cm]
      \centering\safeimg[width=\textwidth,height=3.0cm,keepaspectratio]{credit.pdf}\\[0.05cm]
      \raggedright\footnotesize Free + outstanding credits $= W$. Little's law: residence $\approx W/\mu$.
    \end{column}
  \end{columns}
  \mlsyscredit{Board: Dr Ian Mitchell, CC BY-SA 2.5, Wikimedia Commons. Multi-Agent Coordination, \S18.7.4 (hypothetical rates). Derived in images/calcs.py.}
  \note{[1.5 min] [OPTIONAL] A rate mismatch inside a pipeline of agents. A producer lists files and emits candidate diffs at 5 a second; a reviewer writes a 500-token critique at 40 tokens a second, so it serves 0.08 a second. With no bound, after 10 minutes 2,952 diffs wait, and the newest waits 10.2 hours, long after the repository it was built against has moved: Lecture 4's stale context, now between agents. Per-agent token budgets do not help, because each diff is cheap. A credit window admits a new item only when the consumer frees a credit. Free plus outstanding credits always equal the window, so congestion slows the producer instead of growing a queue. With a saturated window of 8, Lecture 5's Little's law gives mean residence of 8 over 0.08, 100 seconds. A kanban column with a limit is the same rule, used in factories; Lecture 15's collection credits are another. ENGAGE: what window keeps mean residence under a minute? (Expected: 4, derived.) WARN: cancellation, status, and settlement messages must travel outside the data credits, or a full window blocks the signals that free it. FLEX: [OPTIONAL].}
\end{frame}

\begin{frame}{Cancellation closes admission; admitted work must still settle}
  \centering
  \mlsyssteps[width=0.94\textwidth,height=5.7cm,keepaspectratio]{cancel_step}{3}
  \mlsyscredit{Multi-Agent Coordination, \S18.7.1, \S18.8.2, and Tables 7 and 8.}
  \note{[2.5 min] [CORE] Subtree cancellation. Step 1: every grant is bound to the generation counter of each ancestor in its lineage, here generation 7 of the migration task. The interface worker's test job was admitted at the gateway before anything happened. Step 2: the supervisor cancels the task, which advances its generation to 8. The gateway checks every ancestor's generation at admission, so the caller's next tool call is rejected even if the worker has not yet heard about the cancel. One counter invalidates the whole subtree. Step 3: the test job admitted before the cancel keeps its operation identity, owner, reservation, and record until its outcome and charges settle; Lecture 10's settlement rule now spans a tree, and Lecture 12's cleanup graph orders the teardown under a separate administrative credential. The fallacy as true or false: release the parent's reservations as soon as cancellation is requested. False; in the chapter's cohort, failed and canceled work was \$6 of \$18 spent. WARN: a remote gateway that cannot reach the authoritative state must fail closed, or declare a bounded stale-policy window. FLEX: [CORE].}
\end{frame}

\begin{frame}{Each hop narrows authority; allowances are reserved}
  \centering
  \mlsyssteps[width=0.94\textwidth,height=5.2cm,keepaspectratio]{atten_step}{2}\\[0.05cm]
  \small Child grant $\subseteq$ parent grant. Permissions can be shared; money and tokens are partitioned.
  \mlsyscredit{Multi-Agent Coordination, \S18.8.1 and Exercise 7 (allowances); Saltzer and Schroeder 1975; Birgisson et al. 2014.}
  \note{[2.5 min] [CORE] Attenuated delegation. Step 1: each hop passes authority no broader than the delegator's own: the root can read, edit callers, and publish; the caller worker can only edit callers; its inspection child can only read one file. This is Lecture 8's capability grants and attenuation, applied hop by hop, so authority travels in the reference a child holds. Chained-caveat tokens let a holder add restrictions without contacting the issuer. Budgets differ from permissions. Reading does not use up read access, but each dollar can be spent once, so allowances are partitioned: the inspection child's \$5 comes out of the caller's \$13, never from a fresh pool. Step 2: the inspection child names the root as its parent. A runtime that resolved authority from that claimed name would hand it publication rights, Lecture 7's confused deputy built by the delegation tree. The runtime assigns the child's identity and grant; the model has no say. ENGAGE: a worker needs to publish, which its grant forbids. Where does the request go? (Expected: to the task owner, bound to the exact candidate hash under Lecture 9's action-bound approval, and rechecked against lineage at dispatch.) FLEX: [CORE].}
\end{frame}

\begin{frame}{Estimate: three children each ask for \$20}
  \mlsysactivity{Estimate}{2 min}
  \begin{columns}[c]
    \begin{column}{0.44\textwidth}
      \begin{mlsyscard}{mlsysaccent}
        \small\textbf{Given (Exercise 7)}\\[0.05cm]
        Root remaining: \$50\\
        Finalization reserve: \$5\\
        Root's own open work: \$6\\
        Three children ask \$20 each: ``less than \$50''
      \end{mlsyscard}
    \end{column}
    \begin{column}{0.53\textwidth}
      \small
      \uncover<2->{Available to children: $50-5-6 = \mathbf{\$39}$\\[0.1cm]
      Copied from the balance: \$60 of liability\\[0.1cm]}
      \uncover<3->{Atomic reservation in order: first \textbf{admitted} (\$19 left); second and third \textbf{refused}\\[0.1cm]
      \textcolor{mlsysaccent}{\textbf{All three fit at \$13 each.}}}
    \end{column}
  \end{columns}
  \mlsyscredit{Multi-Agent Coordination, Exercise 7(a) (remediation agent; hypothetical amounts). Derived in images/calcs.py.}
  \note{[2 min] [CORE] The remediation agent from Lecture 1 restores a degraded service and proposes three children. ENGAGE: two minutes; what can the runtime admit? Step 2: the children may draw only on \$50 minus the \$5 finalization reserve minus \$6 already reserved for the root, \$39. Copying each child's \$20 from the parent's balance creates \$60 of liability against \$39 of money, which is how concurrent siblings overdraw a task. Step 3: atomic reservation in dispatch order admits the first child and leaves \$19, so the second and third are refused; an equal split admits all three at \$13 each. A child that subdivides further reserves out of its own allowance. WARN: students who answer \$16.67 forgot the reserve and the root's own work. LINK: Lecture 19 takes this account to its conclusion, budget conservation, including when a canceled child's unspent allowance may return. FLEX: [CORE].}
\end{frame}

% =============================================================================
% CLOSE
% =============================================================================
\begin{frame}{Retrieval: which control prevents each failure?}
  \mlsysactivity{Retrieval}{2 min}
  \begin{columns}[c]
    \begin{column}{0.52\textwidth}
      \centering\safeimg[width=\textwidth,height=4.8cm,keepaspectratio]{mast.pdf}
    \end{column}
    \begin{column}{0.45\textwidth}
      {\footnotesize\textcolor{figviolet}{\textbf{design}} \textcolor{figorange}{\textbf{misalignment}} \textcolor{figred}{\textbf{verification}}}\\[0.15cm]
      \only<1|handout:0>{\small Close your notes. For three bars, name the control from today that would catch it.}%
      \only<2->{\footnotesize
      \begin{tabular}{@{}P{2.3cm}P{3.6cm}@{}}\toprule
        \textbf{Failure} & \textbf{Control}\\\midrule
        Disobey task spec & typed envelope\\[0.03cm]
        Unaware of termination & completion contract, budgets\\[0.03cm]
        No or incorrect verification & protected check of the exact candidate\\\bottomrule
      \end{tabular}}
    \end{column}
  \end{columns}
  \mlsyscredit{Cemri et al., ``Why Do Multi-Agent LLM Systems Fail?'', NeurIPS 2025 Datasets and Benchmarks, arXiv 2503.13657v3, Figure 1 (7 of 14 failure modes; 1,642 traces, 7 frameworks).}
  \note{[2 min] [CORE] Real data on real multi-agent systems. Cemri and coauthors annotated 1,642 execution traces from seven open-source multi-agent frameworks and found 14 failure modes in three categories: system design 44.2 percent, inter-agent misalignment 32.3 percent, task verification 23.5 percent. The bars are the seven most common. Retrieval: 90 seconds, close notes, pick three bars and name today's control. Step 2: disobeying the task specification is a prose handoff without a typed envelope; not knowing when to stop needs a completion contract and admission bounds; missing or incorrect verification is what a protected check on the exact merged candidate supplies, where reviewer agreement cannot. The same paper reports that a verification step alone raised one framework's task success by 15.6 percent. ENGAGE: which category does none of today's controls address? (Expected: reasoning-action mismatch inside one agent; that is a model problem, the subject of Part V.) WARN: categories were labeled by people and then an LLM judge; the shares describe these frameworks and benchmarks. FLEX: [CORE].}
\end{frame}

\begin{frame}{Delegation bought deadline coverage at 1.5 times the cost}
  \begin{columns}[c]
    \begin{column}{0.56\textwidth}
      \footnotesize
      \begin{tabular}{@{}lrr@{}}\toprule
        \textbf{100 migrations, \$20 cap} & \textbf{Single} & \textbf{Delegated}\\\midrule
        Accepted by 600 s & 80 & 80\\
        Accepted by 300 s (need 75) & 70 & \textbf{78}\\
        Total spend & \$12 & \$18\\
        Failed and canceled spend & \$3 & \$6\\
        Cost per accepted task & \$0.150 & \$0.225\\\bottomrule
      \end{tabular}\\[0.2cm]
      \small \$6 more buys 8 more on time: \textbf{\$0.75 each}.
    \end{column}
    \begin{column}{0.41\textwidth}
      \centering\safeimg[width=\textwidth,height=2.5cm,keepaspectratio]{research_arch.png}\\[0.1cm]
      \small\raggedright Research: breadth-first, few dependencies, 21 workers can pay.\\[0.05cm]
      Repair: coupled changes; 6 at most, and only for a deadline.
    \end{column}
  \end{columns}
  \mlsyscredit{Multi-Agent Coordination, Napkin Math 18.7 and Table 9 (illustrative cohort); Lighthouse 18.1. Figure: Hadfield et al., 2025. Derived in images/calcs.py.}
  \note{[2.5 min] [CORE] Return to the hook. The matched trial of Napkin Math 18.7: 100 migrations, the same tools, concurrency, and \$20 ceiling for both designs. Both accept 80. Only the delegated design meets the screen of 75 accepted by the 300-second deadline: 78 against 70. It spends \$18 against \$12, so cost per accepted task rises 1.5 times, \$0.225 against \$0.150, and each extra on-time acceptance costs \$0.75 (derived). Delegate if an on-time migration is worth more than \$0.75 to the service; keep the single controller if the objective is cost per accepted task. A third of the delegated spend went to failed and canceled work, the first place to cut the tax. Now the research system: its queries are breadth-first and its search paths rarely reconcile, the regime where the scalability law allowed 21 workers. Its builders say coding tasks parallelize less. Our migration's coupled interface and callers allow 6 at most. ENGAGE: which of today's numbers would you measure first before deploying delegation? (Expected: the single controller's acceptance on your own tasks; above 26.7 percent at a 3.75 times multiple, cost cannot win.) FLEX: [CORE].}
\end{frame}

\begin{frame}{Takeaways}
  \thumbline{thumb_screen.pdf}{The serial fraction caps speedup and overhead sets the useful fleet; at 3.75 times the tokens, cost wins only below 26.7\% single-agent success.}\\[0.12cm]
  \thumbline{thumb_structures_step3.pdf}{Dependencies, messages, and ownership are separate graphs; the runtime admits each child and binds it with a typed envelope.}\\[0.12cm]
  \thumbline{thumb_merge_step3.pdf}{A clean merge does not make the change valid, so revalidate, check the exact candidate, then compare and swap.}\\[0.12cm]
  \thumbline{thumb_floor_step2.pdf}{Correlated reviewers hit a floor (15.0\% here); a different kind of check adds evidence that more reviewers cannot.}\\[0.12cm]
  \thumbline{thumb_cancel_step3.pdf}{Cancellation closes admission; grants narrow per hop, allowances are reserved, and admitted work settles.}
  \note{[1 min] [CORE] One line per segment. Ask students which slide taught each line. If short on time, read the first and the third.}
\end{frame}

\begin{frame}{Open questions}
  \begin{columns}[c]
    \begin{column}{0.62\textwidth}
      \small
      \begin{itemize}\setlength{\itemsep}{0.12cm}
        \item Reviewer correlation must be measured on known-invalid candidates. Where does a team get enough of them for its own code?
        \item Agents can talk in shared chats or pass typed receipts. When is the richer channel worth its quadratic cost?
        \item As models improve, does a single controller with parallel tools make most delegation unnecessary?
      \end{itemize}
    \end{column}
    \begin{column}{0.35\textwidth}
      \centering\safeimg[width=\textwidth,height=3.4cm,keepaspectratio]{shuttle_gpc.jpg}\\[0.04cm]
      {\footnotesize NASA paid a second team for independent code}
    \end{column}
  \end{columns}
  \mlsyscredit{Photo: NASA, public domain, Wikimedia Commons.}
  \note{[1 min] [OPTIONAL] No single answer. Pick one for discussion or assign one as a short written response. The first follows from Napkin Math 18.4: the floor depends on a correlation nobody can read off a model card, and seeding known-invalid candidates is itself work. The second weighs the group chat's $N(N-1)$ deliveries against what a typed receipt throws away. The third comes from the research system's own observation that token usage explained 80 percent of performance variance on one browsing benchmark: if delegation mainly buys more tokens, a single controller with a larger budget may get the same gain. The photo: the Shuttle's builders paid for a second software team to get independence; reviewer agents get none by default.}
\end{frame}

\begin{frame}{Read and explore}
  \small
  \begin{itemize}\setlength{\itemsep}{0cm}
    \item \textbf{Read:} Multi-Agent Coordination, \S18.2 to \S18.10; Checkpoints 18.1 to 18.3; Napkin Math 18.1, 18.4, 18.6.
    \item Hadfield et al., ``How we built our multi-agent research system,'' 2025: a production system and its token bill.
    \item Knight and Leveson, ``An Experimental Evaluation of the Assumption of Independence in Multiversion Programming,'' 1986.
    \item Gunther, the universal scalability law, 2008.
    \item Dean and Barroso, ``The Tail at Scale,'' 2013: fan-out, stragglers, and hedged requests.
    \item Cemri et al., ``Why Do Multi-Agent LLM Systems Fail?'', 2025: 1,642 failure traces.
  \end{itemize}
  \note{[0.5 min] [CORE] The chapter first, with its three checkpoints; work the screening, reviewer, and break-even notebooks by hand. The research-system post is the hook's source. Knight and Leveson is the experiment that broke the independence assumption for software. Gunther derives the scalability law and separates contention from coherency. Dean and Barroso is the classic on fan-out tails. Cemri et al. is the field catalog of how multi-agent systems fail. FLEX: add the agent framework your course uses, and have students find its fan-out and cancellation controls.}
\end{frame}

\begin{frame}[c]
  \begin{columns}[c]
    \begin{column}{0.34\textwidth}
      \centering\safeimg[width=\textwidth,height=1.6cm,keepaspectratio]{thumb_next_bars.pdf}
    \end{column}
    \begin{column}{0.62\textwidth}
      {\small\textcolor{lightgray}{NEXT LECTURE \textbar{} AGENT ECONOMICS}}\\[0.5cm]
      {\Large\bfseries Why can a model with a tenth of the token price cost more per finished task?}
    \end{column}
  \end{columns}
  \note{[0.5 min] [CORE] Today priced one decision, whether to split, by charging failed and canceled work to the tasks it served. Lecture 19 makes that account general for one agent or many: cost per accepted task over a whole cohort, prefix reuse, retries, faster inference, model routing, and capacity. It opens with a model whose tokens cost a tenth as much and asks why switching can still raise the price of every repair that ships.}
\end{frame}

\end{document}
